\documentclass[11pt,letterpaper]{article}
\usepackage[T1]{fontenc}
\usepackage[utf8]{inputenc}
\usepackage{lmodern}
\usepackage[margin=1in,headheight=15pt]{geometry}
\usepackage{amsmath,amssymb,amsthm,mathtools}
\usepackage{microtype,booktabs,longtable,tabularx,array,enumitem}
\usepackage{xcolor,fancyhdr,needspace}
\definecolor{Ink}{HTML}{243D4A}
\definecolor{Accent}{HTML}{41666A}
\usepackage[colorlinks=true,linkcolor=Ink,citecolor=Accent,urlcolor=Accent]{hyperref}
\usepackage{xurl}
\usepackage[nameinlink,noabbrev]{cleveref}
\hypersetup{pdftitle={Spectral Descent and Exact Ramification in Surcomplex Matrix Calculus},pdfauthor={OpenAI, research assistance for Vladimir Reshetnikov},pdfsubject={Arbitrary-rank Hahn fields, exponent-preserving spectral resolution, determinantal ramification, and primitive traces}}
\pagestyle{fancy}\fancyhf{}
\fancyhead[L]{\small Spectral descent and exact ramification}
\fancyhead[R]{\small\thepage}
\renewcommand{\headrulewidth}{0.3pt}
\setlength{\parskip}{0.22em}
\setlength{\emergencystretch}{3em}
\widowpenalty=10000\clubpenalty=10000\displaywidowpenalty=10000
\setlist{itemsep=0.2em,topsep=0.3em}
\setcounter{tocdepth}{1}
\numberwithin{equation}{section}
\newtheorem{theorem}{Theorem}[section]
\newtheorem{lemma}[theorem]{Lemma}
\newtheorem{proposition}[theorem]{Proposition}
\newtheorem{corollary}[theorem]{Corollary}
\theoremstyle{definition}
\newtheorem{definition}[theorem]{Definition}
\newtheorem{example}[theorem]{Example}
\theoremstyle{remark}
\newtheorem{remark}[theorem]{Remark}
\newtheorem{warning}[theorem]{Warning}
\newcommand{\R}{\mathbb R}
\newcommand{\C}{\mathbb C}
\newcommand{\Q}{\mathbb Q}
\newcommand{\Z}{\mathbb Z}
\newcommand{\N}{\mathbb N}
\newcommand{\No}{\mathbf{No}}
\newcommand{\SC}{\mathbf{No}[i]}
\newcommand{\OO}{\mathcal O}
\newcommand{\mm}{\mathfrak m}
\newcommand{\Mat}{\operatorname{Mat}}
\newcommand{\diag}{\operatorname{diag}}
\newcommand{\tr}{\operatorname{tr}}
\newcommand{\rank}{\operatorname{rank}}
\newcommand{\supp}{\operatorname{supp}}
\newcommand{\res}{\operatorname{res}}
\newcommand{\Res}{\operatorname{Res}}
\newcommand{\GL}{\operatorname{GL}}
\newcommand{\Gal}{\operatorname{Gal}}
\newcommand{\Hom}{\operatorname{Hom}}
\newcommand{\Span}{\operatorname{span}}
\newcommand{\st}{\operatorname{st}}
\newcommand{\ord}{\operatorname{ord}}
\newcommand{\lcm}{\operatorname{lcm}}
\newcommand{\abs}[1]{\lvert#1\rvert}
\newcommand{\ang}[1]{\langle#1\rangle}
\newcommand{\KG}{K_{\Gamma}}
\newcommand{\FG}{F_{\Gamma}}
\newcommand{\Commit}{\nolinkurl{aa846271b4dcae2c055b216126a87210292ec19b}}
\newcolumntype{Y}{>{\raggedright\arraybackslash}X}
\newcolumntype{L}[1]{>{\raggedright\arraybackslash}p{#1}}

\begin{document}
\hypersetup{pageanchor=false}
\begin{titlepage}
\centering
\vspace*{0.35in}
{\large\scshape A research article with complete mathematical proofs\par}
\vspace{0.45in}
{\LARGE\bfseries\color{Ink} Spectral Descent and Exact Ramification\\[0.25em]
in Surcomplex Matrix Calculus\par}
\vspace{0.35in}
{\Large Arbitrary-rank Hahn fields, support-certified diagonalization,\\[0.15em]
and primitive traces of positive roots\par}
\vspace{0.4in}
{\large Prepared for Vladimir Reshetnikov\par}
{\large September 21, 2026\par}
\vspace{0.4in}
\begin{minipage}{0.94\textwidth}
\begin{abstract}
\noindent
Let $\Gamma$ be any set-sized ordered abelian group, not necessarily divisible,
and let $\KG=\C((t^\Gamma))$. We give an exponent-preserving spectral construction
for Hermitian and normal matrices over $\KG$. Its support control uses a finite
splitting tree and strongly summable Hahn substitution, not a convergence
argument or passage to a real closure. Singular-value decompositions also stay
in the original exponent group.

The main arithmetic consequence is an exact determinantal ramification law.
For a positive semidefinite Hermitian matrix $A$ of rank $r$, put
$\Delta_0=0$ and $\Delta_k=\min_{I,J}v(\det A_{I,J})$, with $k$-element index
sets. For every integer $q\ge2$, the field generated by the entries of the
principal root $A^{1/q}$ is
\[
 \C\bigl((t^{\Gamma+\sum_{k=1}^{r}\Z\Delta_k/q})\bigr).
\]
Its degree over $\KG$ is the order of the subgroup generated by the classes
$\Delta_k+q\Gamma$ in $\Gamma/q\Gamma$. More unexpectedly,
$\tr(A^{1/q})$ alone is a primitive generator of this entire field. Positivity
prevents cancellation between the relevant monomial cosets. We also prove
finite-witness simultaneous diagonalization, support-compatible transport,
and a local valuation estimate for the splitting projectors.

For surcomplex matrices, the results descend to the subgroup generated by the
actual input supports. Thus eigenvectors need no new fractional scales, while
positive matrix roots have an exactly measurable, and sometimes maximal,
ramification cost.
\end{abstract}
\end{minipage}
\vfill
\begin{minipage}{0.94\textwidth}\small
\textbf{Research status.} The article selects the ``new theorems'' alternative
in the request, not a claimed solution of a named published open problem.
Classical spectral theory, older power-series spectral theorems, Hahn support
lemmas, and Kummer theory are acknowledged as predecessors. The precise
arbitrary-rank determinantal/primitive-trace package is a candidate original
contribution; priority is not certified. The proofs are not independently
refereed or proof-assistant verified. Executed finite symbolic checks accompany
the article and have a narrower scope than its theorems.
\end{minipage}
\end{titlepage}
\hypersetup{pageanchor=true}
\pagenumbering{roman}
\begingroup\small\setlength{\parskip}{0pt}\tableofcontents\endgroup
\clearpage\pagenumbering{arabic}

\section{The problem and the main results}
\label{sec:intro}

\subsection{Why existence over the full surreal field is not the right endpoint}

An existence theorem over a very large ambient field can conceal substantial
arithmetic. A matrix with entries using only integer powers of an infinitesimal
may be diagonalizable without introducing its square root, even when a
positive square root of the matrix itself genuinely requires that new scale.
Passing immediately to a divisible exponent group erases this distinction.

This article retains the exact group of exponents. We work with
\begin{equation}\label{eq:fields}
 \FG=\R((t^\Gamma)),\qquad \KG=\C((t^\Gamma))=\FG[i],
\end{equation}
where $\Gamma$ is an arbitrary ordered abelian group that is a set. A positive
exponent denotes a small monomial. Complex conjugation fixes every $t^\gamma$.
No divisibility, finite-rank, countability, or Archimedean assumption is made
on $\Gamma$.

The resulting question is not merely whether eigenvalues exist in $\SC$.
It is whether a matrix's spectral data can be constructed inside the Hahn
field generated by its original scales, and, when an operation does require
new scales, how to measure the exact extension from the original entries.
The answer separates spectral decomposition from positive root extraction.

\subsection{Three main theorems}

The following statements collect the principal results; their proofs occupy
\cref{sec:support,sec:split,sec:spectral,sec:minors,sec:roots,sec:trace}.
The notation $K(B)$ means the field obtained from $K$ by adjoining all entries
of a finite matrix $B$.

\begin{theorem}[Spectral descent]\label{thm:headline-descent}
Every Hermitian matrix over $\KG$ is unitarily diagonalizable over $\KG$,
with all eigenvalues in $\FG$. Every normal matrix over $\KG$ is unitarily
diagonalizable over $\KG$, and every rectangular matrix has a singular-value
decomposition over these same fields.

If the supports of the input entries lie in a subgroup $H\subseteq\Gamma$,
all the decompositions may be chosen over $K_H$. In the Hermitian case,
a support-certified construction has at most $n-1$ nontrivial splitting
nodes for an $n\times n$ matrix.
\end{theorem}

\begin{theorem}[Determinantal ramification law]\label{thm:headline-root}
Let $A\in\Mat_n(\KG)$ be Hermitian and positive semidefinite, of rank $r$.
For $1\le k\le r$, define
\begin{equation}\label{eq:delta-intro}
 \Delta_k(A)=\min_{\substack{I,J\subseteq\{1,\ldots,n\}\\|I|=|J|=k}}
          v(\det A_{I,J}),\qquad \Delta_0(A)=0.
\end{equation}
Set
\begin{equation}\label{eq:root-group-intro}
 \Gamma_{A,q}=\Gamma+\sum_{k=1}^{r}\Z\frac{\Delta_k(A)}q
       \ \subseteq\ \Gamma\otimes_\Z\Q,
 \qquad
 R_q(A)=\ang{\Delta_1(A)+q\Gamma,\ldots,\Delta_r(A)+q\Gamma}.
\end{equation}
For every integer $q\ge2$, the principal positive semidefinite root satisfies
\begin{align}
 \KG(A^{1/q})&=K_{\Gamma_{A,q}},\label{eq:field-main}\\
 [\KG(A^{1/q}):\KG]&=|R_q(A)|\le q^r.\label{eq:degree-main}
\end{align}
In particular,
\begin{equation}\label{eq:criterion-main}
 A^{1/q}\in\Mat_n(\KG)
 \quad\Longleftrightarrow\quad
 \Delta_k(A)\in q\Gamma\quad(1\le k\le r).
\end{equation}
The extension is finite, abelian, and totally ramified, with Galois group
isomorphic to $\Hom(R_q(A),\mu_q)$, where $\mu_q\subset\C^\times$ is the group
of $q$th roots of unity.
\end{theorem}

\begin{theorem}[Primitive trace]\label{thm:headline-trace}
Under the hypotheses of \cref{thm:headline-root},
\begin{equation}\label{eq:trace-main}
 \KG\bigl(\tr(A^{1/q})\bigr)=\KG(A^{1/q})=K_{\Gamma_{A,q}}.
\end{equation}
More generally, if $\lambda_1,\ldots,\lambda_r$ are the positive eigenvalues
of $A$, then every sum $\sum_{j=1}^{r}a_j\lambda_j^{1/q}$ with
$a_j\in\FG$ and $a_j>0$ is a primitive generator of this same field.
\end{theorem}

The last statement is not an assertion that an arbitrary linear combination
of radicals is primitive. Strictly positive coefficients are what force every
necessary coset to remain visible. Nor is \cref{eq:degree-main} simply a bound:
it computes the exact degree using finitely many minor valuations, without
first computing an eigenbasis or the eigenvalues themselves.

\subsection{What these statements add to the repository}

The repository snapshot inspected for this article is
\begin{center}\small\Commit.\end{center}
Its spectral report already develops the SVD over real closed fields,
determinantal singular scales, positive Cauchy--Binet identities, and
perturbation theory. It works in divisible Hahn workspaces and explicitly
separates pointwise diagonalization from a support-controlled construction
\cite{RepoSpectral}. Those existing contributions are not claimed again as
new here.

The change is to \emph{nondivisible} workspaces, followed by an exact
field-of-definition calculation. In particular, the determinantal invariant
already useful for singular scales now controls a finite radical extension,
and the trace of the principal root detects all of that extension. The
support argument explains why no hidden exponent enlargement is needed in
the spectral step preceding that calculation.

The audit was targeted: the repository map, its main README, the full spectral
README, and selected neighboring material were inspected. It was not a
line-by-line review of every archived source manuscript. Absence from a search
result is not used as proof of absence from the collection.

\subsection{Published predecessors and the precise novelty boundary}
\label{sec:prior}

There is substantial prior spectral theory over power-series fields.
Adkins established a spectral theorem for normal matrices over hermitian
discrete valuation rings, including formal and convergent one-variable power
series \cite{Adkins}. Keller and Ochsenius proved orthogonal diagonalization
over iterated real Laurent-series fields and a particular infinite-rank Hahn
field, using recursive block splitting \cite{KO}. Thus neither ``a spectral
theorem beyond $\R$'' nor ``no fractional exponents in a Laurent-series
eigenbasis'' is a new claim.

Parusi\'nski and Rond prove multiparameter formal spectral and singular-value
results under normal-crossing discriminant hypotheses \cite{PR}. Their work
contains the formal separated-block lifting ingredient that underlies many
such constructions. The 2026 preprint of Dai, Liang, Lu, and Zhi gives a
splitting-and-projector criterion over multivariable formal power-series rings
and an algorithmic treatment over regular local rings \cite{DLLZ}. Our Hahn
field is a different ring: a totally ordered exponent group permits monomial
division that a multivariable power-series ring need not permit. We do not
claim to remove the obstructions in that latter category.

The support foundations below are standard Hahn--Neumann and word-order
ideas \cite{Neumann,Higman}; the finite extension calculation is an elementary
monomial form of Kummer theory. What is offered as a \emph{candidate original
package} is the arbitrary-rank input-support descent together with
\cref{eq:field-main,eq:degree-main,eq:trace-main}, their finite support
certificates, and their exact interpretation inside $\SC$. I did not locate
this precise determinantal/primitive-trace formulation in the sources
reviewed. This is a qualified literature assessment, not a certification of
priority or a claim that each lemma is new.

\section{Hahn fields, positivity, and the meaning of support control}
\label{sec:hahn}

\subsection{The coefficient and exponent conventions}

An element of $k((t^\Gamma))$, for $k=\R$ or $\C$, is a formal expression
\[
 f=\sum_{\gamma\in S}f_\gamma t^\gamma
\]
with $S\subseteq\Gamma$ well ordered. Its support is the set of exponents with
nonzero coefficients, and
\[
 v(f)=\min\supp(f)\quad(f\ne0),\qquad v(0)=+\infty.
\]
Addition is coefficientwise and multiplication is convolution. The Hahn
support rules ensure that all coefficients of a product are finite sums.
The order on $\FG$ is the leading-coefficient order: a nonzero series is
positive exactly when its leading real coefficient is positive. In
particular $t^\gamma>0$ for every $\gamma$, including negative exponents.

Put
\[
 \OO=\{f:v(f)\ge0\},\qquad \mm=\{f:v(f)>0\}.
\]
The coefficient at $0$ gives the residue map $\OO\to\C$, or $\OO\cap\FG\to\R$.
The star operation on matrices is conjugate transpose. We use
\[
 v(x)=\min_jv(x_j),\qquad v(A)=\min_{i,j}v(A_{ij}).
\]
A matrix is Hermitian when $A^*=A$, normal when $A^*A=AA^*$, and unitary when
$U^*U=UU^*=I$. A Hermitian $A$ is positive semidefinite when
$x^*Ax\ge0$ for all $x\in\KG^n$.

\subsection{Strong summation, not a limit convention}

\begin{definition}[Strongly summable family]
A set-indexed family $(f_i)_{i\in I}$ of Hahn series is strongly summable if
$\bigcup_i\supp(f_i)$ is well ordered and, for each exponent $\gamma$, only
finitely many $i$ have a nonzero coefficient at $\gamma$. Its sum is the
coefficientwise finite sum.
\end{definition}

This definition licenses arbitrary regrouping of a jointly summable family.
It does not declare a sequence of partial sums to converge. In particular,
the phrase ``formal power series evaluated at infinitesimals'' below always
comes with a support proof.

\begin{example}[A rank-two warning]\label{ex:nonconvergence}
Let $\Gamma=\Z^2$ with lexicographic order and let
$\delta=(0,1)$, $\Omega=(1,0)$. Then $0<n\delta<\Omega$ for every positive
integer $n$. The Hahn identity
\[
 \sum_{n\ge0}t^{n\delta}=(1-t^\delta)^{-1}
\]
is valid, but the sequence of partial sums does not converge to the right-hand
side in the valuation topology: the tail has valuation $(N+1)\delta<\Omega$
for every $N$. Thus a bound $v(\text{remainder}_N)\ge N\delta$ is not by itself
a convergence proof in arbitrary rank. It will never be used as one here.
\end{example}

\subsection{Norms exist even when positive square roots do not all exist}

\begin{lemma}[No leading cancellation in sums of squares]\label{lem:norm}
For every nonzero $x\in\KG^n$,
\begin{equation}\label{eq:norm-val}
 x^*x>0,\qquad v(x^*x)=2v(x).
\end{equation}
Moreover $x^*x$ has a unique positive square root in $\FG$.
\end{lemma}

\begin{proof}
Let $\alpha=v(x)$ and write $x=t^\alpha y$, where all entries of $y$ are
integral and at least one has a nonzero constant coefficient. The constant
coefficient of $y^*y$ is $c=\sum_j|\res(y_j)|^2>0$. Hence
\[
 x^*x=c\,t^{2\alpha}(1+h),\qquad h\in\mm\cap\FG.
\]
The formal binomial series $(1+h)^{1/2}$ is a Hahn series by
\cref{lem:evaluate} below. Therefore
$\sqrt c\,t^\alpha(1+h)^{1/2}\in\FG$ is the required positive square root.
Uniqueness follows from $(u-v)(u+v)=0$ in an ordered field.
\end{proof}

This proof does not assume that $\FG$ is real closed. For example, if
$\gamma\notin2\Gamma$, the positive element $t^\gamma$ has no square root in
$\FG$, since the valuation of a square belongs to $2\Gamma$. Such an element
cannot be a nonzero sum of squares. This distinction is exactly why an SVD
can stay in the base field while an arbitrary positive matrix square root
need not.

\begin{lemma}[Integral unitary matrices]\label{lem:unitary}
Every entry of a unitary matrix belongs to $\OO$. For every vector $x$ and
matrix $B$ of compatible dimensions,
\[
 v(Ux)=v(x),\qquad v(UB)=v(B),\qquad v(BU)=v(B).
\]
\end{lemma}
\begin{proof}
Each column of $U$ has squared norm $1$, so \cref{eq:norm-val} gives minimum
column valuation $0$. Thus every entry is integral. Integral multiplication
cannot lower minimum valuation. Applying this both to $U$ and its integral
inverse $U^*$ proves the three equalities.
\end{proof}

\section{A self-contained support engine}
\label{sec:support}

\subsection{The word lemma and the Hahn--Neumann lemma}

We record proofs because finiteness \emph{across all word lengths}, not merely
at a fixed degree, is indispensable in higher rank. The combinatorial method
is the standard minimal-bad-sequence argument associated with Higman's lemma
\cite{Higman}; the resulting support statement is the Hahn--Neumann lemma
\cite{Neumann}.

\begin{lemma}[Ordered words]\label{lem:words}
Let $S$ be a well-ordered set. Order finite words in $S$ by subsequence
embedding with coordinatewise increase. Every infinite sequence of words has
a pair $i<j$ such that the $i$th word embeds in the $j$th.
\end{lemma}
\begin{proof}
Suppose a bad sequence exists. Choose it successively so that at each position
the word has the least possible length among words allowing a bad
continuation. No chosen word is empty. Write $w_i=u_i a_i$ with last letter
$a_i$. A sequence in a well order has an infinite nondecreasing subsequence;
choose $a_{i_0}\le a_{i_1}\le\cdots$.

Then $w_0,\ldots,w_{i_0-1},u_{i_0},u_{i_1},\ldots$ is still bad. Indeed, a
comparison from its initial segment into a $u_{i_j}$ would compare the
corresponding original words. A comparison $u_{i_j}\preceq u_{i_k}$, together
with $a_{i_j}\le a_{i_k}$, would compare $w_{i_j}$ and $w_{i_k}$. Both are
impossible. This contradicts the minimal length at position $i_0$.
\end{proof}

\begin{lemma}[Positive support monoids]\label{lem:neumann}
Let $S\subset\Gamma_{>0}$ be well ordered. The additive monoid
$M=\ang{S}_{\N}$ is well ordered, and each $\gamma\in\Gamma$ is the sum of
only finitely many finite ordered words with letters in $S$.
\end{lemma}
\begin{proof}
If $M$ admitted an infinite strictly decreasing sequence, choose a word
representing each term. \Cref{lem:words} gives an embedding from an earlier
word into a later one. Positivity of unused letters and increase of matched
letters make the earlier sum at most the later sum, a contradiction.

If infinitely many distinct words had sum $\gamma$, apply the same lemma to
those words. Equality of sums forbids any unused positive letter or any
strict increase of a matched letter. The two words must consequently have
the same length and the same letters, contrary to their being distinct.
\end{proof}

\subsection{Substitution of finite-variable formal series}

\begin{lemma}[Hahn evaluation]\label{lem:evaluate}
Let $z_1,\ldots,z_d\in\mm$ and let
$S=\bigcup_{j=1}^d\supp(z_j)$. Every formal series
$f\in\C[[X_1,\ldots,X_d]]$ has a well-defined strongly summed evaluation
$f(z_1,\ldots,z_d)\in\KG$. Its support lies in
$M=\ang{S}_{\N}$; if $f(0)=0$, the support lies in $M\setminus\{0\}$.
Evaluation preserves sums, products, conjugation with real formal variables,
and formal composition whose inner series have zero constant coefficient.
\end{lemma}
\begin{proof}
A finite union of well-ordered subsets of a linear order is well ordered.
Every contribution to a coefficient of $f(z)$ determines an ordered word of
exponents in $S$, together with finitely many choices of one of the $d$
variables at each position. By \cref{lem:neumann}, for a fixed exponent there
are only finitely many such words and labels. Thus all coefficient sums are
finite. Their support is contained in the well-ordered monoid $M$.

For a product or a composition, expanding either side gives the same labeled
words with the same finite contributions at each exponent. Hence regrouping
is legal and the formal identities survive evaluation. Conjugation preserves
supports and the same coefficientwise argument applies.
\end{proof}

Matrix formal series reduce to this lemma by using finitely many scalar
entries as variables. The entries need not commute as matrices; scalar
coefficient calculations keep track of all ordered matrix products.
No bound on the growth of the ordinary complex Taylor coefficients is
required: a coefficient at a fixed Hahn exponent is still a finite sum.

\begin{corollary}[Exponent-preserving inverses and unit roots]\label{cor:unitroots}
If $h\in\mm$, the geometric and binomial series define
$(1+h)^{-1}$ and $(1+h)^a$ for every rational $a$. All their exponents lie in
the group generated by $\supp(h)$. Every nonzero Hahn series has an inverse
without enlarging its exponent group.
\end{corollary}
\begin{proof}
Use the formal geometric and binomial identities and \cref{lem:evaluate}.
For a general series factor $f=c t^\alpha(1+h)$. Its inverse is
$c^{-1}t^{-\alpha}(1+h)^{-1}$, and both $\alpha$ and the exponents of $h$
belong to the original support subgroup.
\end{proof}

\subsection{A valuation estimate for formal substitution}

\begin{lemma}[Formal maps are valuation nonexpanding]\label{lem:lipschitz}
Let $f\in\C[[X_1,\ldots,X_d]]$ and let $x,y\in\mm^d$. Then
\[
 v\bigl(f(x)-f(y)\bigr)\ \ge\ \min_jv(x_j-y_j).
\]
The same entrywise assertion holds for a matrix of formal series.
\end{lemma}
\begin{proof}
There are formal series $g_j\in\C[[X,Y]]$ with
\[
 f(X)-f(Y)=\sum_{j=1}^d(X_j-Y_j)g_j(X,Y).
\]
For a monomial this is a telescoping product identity; collecting it by total
degree gives the formal identity. Evaluate with \cref{lem:evaluate}.
Each $g_j(x,y)$ is integral, so the claimed inequality follows.
\end{proof}

This elementary estimate concerns actual Hahn series after certified
substitution. It does not assert that all positive valuations are integer
multiples of a single scale.

\section{Support-certified separation of spectral clusters}
\label{sec:split}

\subsection{Universal projectors}

Let
\begin{equation}\label{eq:D}
 D=\diag(c_1I_{n_1},\ldots,c_sI_{n_s}),
 \qquad c_j\in\R,\quad c_j\ne c_\ell\ (j\ne\ell),
\end{equation}
and write $Q_j$ for the corresponding constant coordinate projections.
Suppose $E=E^*\in\Mat_n(\mm)$ and $B=D+E$. Let $S$ be the union of the
supports of the entries of $E$ and let $M=\ang{S}_{\N}$.

\begin{theorem}[Separated-block lifting]\label{thm:split}
There are Hermitian projections $P_1,\ldots,P_s$ and a unitary matrix $W$
over $\KG$ such that
\begin{align}
 &P_jP_\ell=0\ (j\ne\ell),\quad \sum_jP_j=I,\quad BP_j=P_jB,
       \quad P_j\equiv Q_j\pmod{\mm},\label{eq:projectors}\\
 &W\equiv I\pmod{\mm},\qquad P_jW=WQ_j,\qquad
 W^*BW\text{ is block diagonal relative to }(Q_j).\label{eq:W}
\end{align}
Every entry of $P_j-Q_j$ and $W-I$ is supported in $M\setminus\{0\}$.
The construction uses universal finite-variable formal series in the entries
of $E$, with ordinary constant coefficients.
\end{theorem}

\begin{proof}
Put $R_0(\zeta)=(\zeta I-D)^{-1}$. Define
\begin{equation}\label{eq:proj-resolvent}
 P_j(E)=\sum_{k\ge0}\Res_{\zeta=c_j}
       \left[R_0(\zeta)\bigl(E R_0(\zeta)\bigr)^k\right].
\end{equation}
For each $k$ the residue is a homogeneous polynomial of degree $k$ in the
entries of $E$. Its coefficients are ordinary complex numbers obtained by
residues of rational functions with poles among the fixed real $c_j$.
Consequently \cref{lem:evaluate} makes \cref{eq:proj-resolvent} a strongly
summed Hahn series with the stated support. The term $k=0$ is $Q_j$.

We justify all formal projector identities without presupposing a spectral
theorem over $\KG$. Parametrize Hermitian matrices $E$ by $n^2$ independent
real coordinates. For sufficiently small \emph{ordinary real} coordinates,
disjoint complex circles around the $c_j$ separate the eigenvalues of the
ordinary complex Hermitian matrix $D+E$. To see that no eigenvalue escapes
these circles, note that if $\lambda$ is farther than $\|E\|$ from every
$c_j$, then $I-(\lambda I-D)^{-1}E$ is invertible by its ordinary geometric
series, so $\lambda I-D-E$ is invertible.

On those ordinary circles the ordinary resolvent has its Neumann expansion.
Its contour integrals are the orthogonal spectral projectors, as follows
immediately by diagonalizing the ordinary Hermitian matrix. Their Taylor
series at $E=0$ are exactly \cref{eq:proj-resolvent}. The identities
$P_j^*=P_j$, $P_j^2=P_j$, $P_jP_\ell=0$, $\sum_jP_j=I$, and $[B,P_j]=0$
therefore hold as finite-variable formal identities: every Taylor coefficient
of their difference is zero. Hahn evaluation transports these identities to
our $E$. This is a proof of universal coefficient identities using ordinary
complex analysis, \emph{not} contour integration in a surreal topology.

Now set
\begin{equation}\label{eq:polar-intertwiner}
 X=\sum_jP_jQ_j,
 \qquad G=X^*X=\sum_jQ_jP_jQ_j,
 \qquad W=XG^{-1/2}.
\end{equation}
Here $X=I+\Mat_n(\mm)$ and $G=I+H$ with $H\in\Mat_n(\mm)$. The series
\begin{equation}\label{eq:Groot}
 G^{-1/2}=\sum_{k\ge0}\binom{-1/2}{k}H^k
\end{equation}
is strongly summable by \cref{lem:evaluate}. It is Hermitian, commutes with
$G$, and squares to $G^{-1}$. Therefore $W^*W=I$, and a square matrix with a
left inverse is invertible, so also $WW^*=I$.

The identities $P_jX=XQ_j$ follow from orthogonality of the projectors.
The matrix $G$, and hence $G^{-1/2}$, commutes with every $Q_j$.
Thus $P_jW=WQ_j$. Since $B$ commutes with $P_j$, its conjugate $W^*BW$
commutes with every $Q_j$, which is precisely block diagonality.
All nonconstant terms in these constructions are formal series without
constant term in the entries of $E$; their supports lie in
$M\setminus\{0\}$.
\end{proof}

\begin{remark}[What is canonical and what is not]
For the fixed residue decomposition $D$ and projections $Q_j$, formulas
\cref{eq:proj-resolvent,eq:polar-intertwiner} specify a particular lift.
They do not choose a preferred ordinary eigenbasis inside a repeated
eigenspace of a later leading coefficient. Those finite-dimensional choices
will be recorded in the splitting tree.
\end{remark}

\subsection{First-order terms and scale-sensitive stability}

Write $E_{j\ell}=Q_jEQ_\ell$. With formal degree in the entries of $E$,
\begin{equation}\label{eq:first-order}
 W=I+K+O_{\mathrm{formal}}(E^2),\qquad
 K_{j\ell}=\frac{E_{j\ell}}{c_\ell-c_j}\quad(j\ne\ell),\qquad K_{jj}=0.
\end{equation}
The matrix $K$ is skew-Hermitian. The formula follows by taking the $k=1$
residues in \cref{eq:proj-resolvent}; $X=I+K+O(E^2)$ gives
$X^*X=I+O(E^2)$. In particular the first-order off-diagonal term of
$W^*(D+E)W$ is $E+[D,K]$, which vanishes between different blocks.
The diagonal block to second degree is
\begin{equation}\label{eq:second-order}
 c_jI+E_{jj}+
   \sum_{\ell\ne j}\frac{E_{j\ell}E_{\ell j}}{c_j-c_\ell}
       +O_{\mathrm{formal}}(E^3).
\end{equation}
These $O_{\mathrm{formal}}$ symbols mean support-certified sums of monomials
of the indicated total degrees, not a topological limiting convention.

\begin{proposition}[Local valuation estimate]\label{prop:local-stability}
Apply the construction with the same $D,Q_j$ to Hermitian
$E,E'\in\Mat_n(\mm)$. Then
\[
 v(P_j(E)-P_j(E'))\ge v(E-E'),\qquad
 v(W(E)-W(E'))\ge v(E-E').
\]
If $A=\mu I+t^\alpha V(D+E)V^*$, where $V$ is constant unitary, and
$v(A'-A)>\alpha$, then the corresponding cluster projectors of $A,A'$ satisfy
\begin{equation}\label{eq:gap-bound}
 v(\Pi_j(A')-\Pi_j(A))\ge v(A'-A)-\alpha.
\end{equation}
\end{proposition}
\begin{proof}
The first pair of inequalities is \cref{lem:lipschitz} applied to the
universal formal entries. For the second assertion use the same $\mu,\alpha,V$
and put $E'=E+t^{-\alpha}V^*(A'-A)V$. Its entries are infinitesimal, and
\cref{lem:unitary} gives $v(E'-E)=v(A'-A)-\alpha$.
\end{proof}

The estimate is sharp already for a two-by-two constant diagonal matrix
perturbed in an off-diagonal entry: \cref{eq:first-order} has a nonzero
linear term. Its role here is to give a support-compatible valuation version
of separated-cluster stability, not to claim a new classical perturbation
inequality.

\section{Finite splitting trees and exact spectral descent}
\label{sec:spectral}

\subsection{The dimension-decreasing construction}

\begin{theorem}[Hermitian Hahn spectral theorem]\label{thm:hermitian}
Let $\Gamma$ be any set-sized ordered abelian group. For every
$A=A^*\in\Mat_n(\KG)$ there are $U\in\Mat_n(\KG)$ and
$\lambda_1,\ldots,\lambda_n\in\FG$ such that
\[
 U^*U=UU^*=I,\qquad U^*AU=\diag(\lambda_1,\ldots,\lambda_n).
\]
For real symmetric input, $U$ can be chosen real orthogonal. Every exponent
occurring in $U$ or in the $\lambda_j$ belongs to the subgroup generated by
the exponents in the input entries.
\end{theorem}

\begin{proof}
We prove the theorem by induction on $n$, using no algebraic closedness or
real closedness assumption on the Hahn fields. The assertion for $n=1$ is
immediate. Set
\begin{equation}\label{eq:trace-subtract}
 \mu=\frac1n\tr(A),\qquad C=A-\mu I.
\end{equation}
If $C=0$, take $U=I$. Otherwise let $\alpha=v(C)$ and form
\begin{equation}\label{eq:normalize}
 t^{-\alpha}C=C_0+E_0,
 \qquad C_0\in\Mat_n(\C),\quad E_0\in\Mat_n(\mm).
\end{equation}
The constant matrix $C_0$ is nonzero, Hermitian, and traceless. It has at least
two distinct ordinary real eigenvalues. Indeed, by the ordinary spectral
theorem a Hermitian matrix with just one eigenvalue is scalar; a traceless
scalar matrix in characteristic zero is zero.

Choose an ordinary constant unitary $V$ with $V^*C_0V=D$ of the form
\cref{eq:D}, with $s\ge2$. Apply \cref{thm:split} to
$D+V^*E_0V$. Then
\[
 W^*V^*AVW=\diag(A_1,\ldots,A_s)
\]
for Hermitian matrices $A_j$ of strictly smaller sizes. By induction choose
unitary $U_j$ diagonalizing each $A_j$. The product
\[
 U=VW\diag(U_1,\ldots,U_s)
\]
diagonalizes $A$, and all diagonal entries are real because the diagonal
matrix is Hermitian.

Let $H$ be the subgroup generated by the input supports. The scalar $\mu$
and matrix $C$ have supports in $H$, and $\alpha\in H$. Thus normalization
only translates supports inside $H$. Constant changes of basis do not add
exponents. \Cref{thm:split} uses sums of positive exponents already in $H$,
so $W$ and the $A_j$ remain in $K_H$. The induction therefore takes place
inside $K_H$ throughout. For real input use real constant eigenbases; every
coefficient of the subsequent projector and polar constructions is real.
\end{proof}

\begin{corollary}[The splitting tree is finite]\label{cor:tree}
The construction in \cref{thm:hermitian} has at most $n-1$ nontrivial splitting
nodes. This bound is independent of the order type of the input supports and
of the rank or cardinality of the exponent group.
\end{corollary}
\begin{proof}
Every internal node has at least two children; the dimension attached to a
node is the sum of its children's positive dimensions. Hence there are at
most $n$ leaves. A finite rooted tree in which every internal node has at
least two children has at most one fewer internal nodes than leaves.
Termination itself follows by induction on the strictly decreasing positive
dimensions along each branch.
\end{proof}

\subsection{What a support certificate actually records}

At each internal node the construction records the scalar series $\mu$,
the exact leading exponent $\alpha$, the ordinary constant spectral data
$(V,D,Q_j)$, and the well-ordered positive support set $S$ of the normalized
perturbation. It then records the inclusion of the lifted projector and
unitary supports in $\ang{S}_{\N}$, followed by the certificates of the smaller
blocks. A scalar leaf records its scalar series and multiplicity.

This is a \emph{finite tree of set-valued support data}. It is not generally
a finite list of monomials. A single node can carry an infinite or
uncountable well-ordered support. Nor is it automatically an effective
algorithm from an arbitrary presentation of a surreal number: exact zero
tests, support access, and extraction of a leading coefficient must be
supplied by the representation. The coefficient fields remain the full $\R$ and $\C$: preservation of a
smaller coefficient field, such as $\Q$, is not asserted. The construction
proves existence and specifies admissible operations; it does not solve the representation or
computability problems discussed in the repository's computer-algebra report
\cite{RepoCAS}.

\begin{example}[A scalar tail beyond every finite truncation]\label{ex:tail}
In the rank-two group of \cref{ex:nonconvergence}, put
\[
 c=\sum_{m\ge1}t^{m\delta},\qquad
 A=cI+t^\Omega\diag(1,-1).
\]
Every coefficient encountered at the scales $m\delta$ is scalar. A procedure
that simply inspects more and more of those coefficients does not reach the
first nonscalar scale after any finite number of such inspections. The
exact trace subtraction gives $A-\tfrac12\tr(A)I=t^\Omega\diag(1,-1)$
immediately. This explains both the usefulness and the exact-operation
requirements of \cref{eq:trace-subtract}.
\end{example}

\subsection{Normal matrices and finite witnesses for commuting families}

\begin{corollary}[Normal spectral descent]\label{cor:normal}
Every normal matrix over $\KG$ is unitarily diagonalizable over $\KG$, without
an exponent-group extension.
\end{corollary}
\begin{proof}
Write $A=H+iJ$, where $H=(A+A^*)/2$ and $J=(A-A^*)/(2i)$ are Hermitian.
Normality is equivalent to $HJ=JH$. Diagonalize $H$ by
\cref{thm:hermitian}. The equation $HJ=JH$ makes the entries of $J$ between
distinct exact eigenspaces of $H$ zero. Diagonalize the Hermitian restriction
of $J$ inside each remaining block. The resulting unitary diagonalizes both
$H$ and $J$, hence $A$. All operations remain in the original support group.
\end{proof}

\begin{lemma}[A finite-dimensional commutant fact]\label{lem:commutant}
If $A$ is normal and $BA=AB$ over $\KG$, then $BA^*=A^*B$.
\end{lemma}
\begin{proof}
Use \cref{cor:normal} to make $A$ diagonal. The entry $b_{jk}$ must vanish
when $\lambda_j\ne\lambda_k$, and these are exactly the pairs for which
$\overline{\lambda_j}\ne\overline{\lambda_k}$. Thus $B$ also commutes with
the diagonal conjugate of $A$.
\end{proof}

\begin{theorem}[Finite-witness simultaneous diagonalization]\label{thm:simultaneous}
Let $\mathcal F\subseteq\Mat_n(\KG)$ be a set of pairwise commuting normal
matrices. There is a simultaneous unitary diagonalizer. It can be constructed
using at most $n-1$ selected members of $\mathcal F$; consequently its support
may be confined to the subgroup generated by those finitely many members'
supports. After selection, one ordinary complex linear combination of the
selected matrices can be chosen whose eigenspaces are exactly the final
joint eigenspaces.
\end{theorem}
\begin{proof}
Begin with the whole space as one block. If every member acts as a scalar
there, stop. Otherwise choose a nonscalar member and decompose the block
into its distinct eigenspaces by \cref{cor:normal}. Every other member and
its adjoint preserve those spaces by \cref{lem:commutant}; hence their
restrictions remain normal. Repeat on any block on which some member is
nonscalar. Every selection strictly increases the number of joint blocks,
which never exceeds $n$. Thus at most $n-1$ selections are needed. The
successive diagonalizers are produced from the selected matrices and
previously constructed bases only, proving the support assertion.

Let $B_1,\ldots,B_m$ be the selected matrices, and let their joint scalar
tuples on the final blocks be $\boldsymbol\lambda^{(a)}\in\KG^m$. These
tuples are pairwise distinct. For each pair $a\ne b$, the condition
\[
 \sum_{j=1}^m c_j(\lambda_j^{(a)}-\lambda_j^{(b)})=0
\]
defines a proper complex-linear subspace of $\C^m$. A finite union of proper
linear subspaces over the infinite field $\C$ does not fill $\C^m$. Choose
$(c_j)$ outside that union. Then $\sum c_jB_j$ has distinct scalar values
on the final blocks. It is normal because the selected matrices commute
with one another and with one another's adjoints.
\end{proof}

For a family of scalar matrices no selected member is needed. The finite
witness statement is an existence theorem for a set-sized family, not a
procedure for searching an arbitrary undecidable family.

\section{Singular values stay in the original scale group}
\label{sec:svd}

\begin{theorem}[SVD over a nondivisible Hahn field]\label{thm:svd}
Every $A\in\Mat_{m\times n}(\KG)$ admits
\[
 A=U\Sigma V^*,
\]
where $U,V$ are unitary over $\KG$ and $\Sigma$ is rectangular diagonal with
positive entries $\sigma_1\ge\cdots\ge\sigma_r>0$ in $\FG$ and zeros
elsewhere. The factors can be chosen in the subgroup generated by the input
supports.
\end{theorem}
\begin{proof}
Diagonalize $A^*A$ by \cref{thm:hermitian}. If $v_j$ is a unit eigenvector
with eigenvalue $\lambda_j$, then
\[
 \lambda_j=v_j^*A^*Av_j=(Av_j)^*(Av_j)\ge0.
\]
When $\lambda_j>0$, \cref{lem:norm} gives
$\sigma_j=\sqrt{\lambda_j}\in\FG$ without any extension of the group.
Define $u_j=Av_j/\sigma_j$. These vectors are orthonormal, since
\[
 u_j^*u_k=\frac{v_j^*A^*Av_k}{\sigma_j\sigma_k}=\delta_{jk}.
\]
If $\lambda_j=0$, \cref{lem:norm} forces $Av_j=0$. Extend the nonzero $u_j$
to an orthonormal basis by the usual finite Gram--Schmidt procedure, selecting
standard basis vectors not yet in the span. Every nonzero residual vector
has a norm in $\FG$ by \cref{lem:norm}. No step requires a square root of
an arbitrary positive scalar. Reordering the finitely many positive singular
values proves the stated form and preserves the support field.
\end{proof}

\begin{corollary}[Gram square roots and polar decomposition]\label{cor:gramroot}
For every rectangular $A$ over $\KG$, its Gram matrix $A^*A$ has its principal
positive square root in $\Mat_n(\KG)$. For square invertible $A$, the polar
decomposition $A=QP$ has $Q$ unitary and $P$ positive definite over $\KG$.
\end{corollary}
\begin{proof}
With the notation of \cref{thm:svd}, put
$P=V\diag(\sigma_1,\ldots,\sigma_r,0,\ldots,0)V^*$.
Then $P^2=A^*A$. In the invertible square case take $Q=AP^{-1}$, or $Q=UV^*$.
\end{proof}

This corollary does not say that every positive matrix is a Gram matrix over
the nondivisible base field. The scalar matrix $(t^\gamma)$ with
$\gamma\notin2\Gamma$ is a counterexample. A precise characterization of
which positive matrices are Gram matrices will follow from the main law.

\section{Determinantal valuations recover the spectral scales}
\label{sec:minors}

\subsection{Invariance without a discrete valuation}

For an $m\times n$ matrix $A$ of rank $r$, define $\Delta_k(A)$ by the same
minimum over all $k\times k$ minors as in \cref{eq:delta-intro}; set
$\Delta_0(A)=0$ and $\Delta_k(A)=+\infty$ for $k>r$.
There are finitely many minors, so no completeness assumption is required
for this minimum.

\begin{lemma}[Integral equivalence preserves minor minima]\label{lem:minor-invariance}
If $P\in\GL_m(\OO)$ and $Q\in\GL_n(\OO)$ have integral inverses, then
\[
 \Delta_k(PAQ)=\Delta_k(A)
\]
for every $k$.
\end{lemma}
\begin{proof}
Cauchy--Binet expresses each $k$-minor of $PA$ as a sum of products of a
$k$-minor of $P$ and a $k$-minor of $A$. Integral minors of $P$ have
nonnegative valuation, so $\Delta_k(PA)\ge\Delta_k(A)$. Applying the same
argument to $P^{-1}$ gives equality. Right multiplication is identical.
\end{proof}

\begin{proposition}[Spectral and singular scales]\label{prop:scales}
Let a normal matrix $A$ have nonzero eigenvalues whose valuations, in
nondecreasing order, are $\gamma_1\le\cdots\le\gamma_r$. Then
\begin{equation}\label{eq:minor-scales}
 \Delta_k(A)=\gamma_1+\cdots+\gamma_k\quad(1\le k\le r).
\end{equation}
For an arbitrary rectangular matrix the same formula holds with
$\gamma_j=v(\sigma_j)$, the singular-value valuations in nondecreasing order.
\end{proposition}
\begin{proof}
Unitary diagonalization and the SVD use integral matrices with integral
inverses by \cref{lem:unitary}. Apply \cref{lem:minor-invariance}. For a
diagonal matrix, the nonzero minors are products of distinct diagonal
entries, and the least valuation of a $k$-fold product is the sum of the
$k$ smallest valuations.
\end{proof}

The singular-scale identity is already a central invariant in the
repository's spectral report \cite{RepoSpectral}. The short proof here makes
explicit that it is valid before imposing divisibility on the group.
For a positive matrix, it yields
\begin{equation}\label{eq:recover-gamma}
 \gamma_k=\Delta_k(A)-\Delta_{k-1}(A),
\end{equation}
which will remove eigenvalues from the ramification criterion.

\subsection{The Gram case and its parity constraint}

\begin{corollary}[Even determinantal scales of Gram matrices]\label{cor:gram-scales}
For $A\in\Mat_{m\times n}(\KG)$ and $1\le k\le\rank(A)$,
\begin{equation}\label{eq:gram-scales}
 \Delta_k(A^*A)=2\Delta_k(A).
\end{equation}
\end{corollary}
\begin{proof}
The nonzero eigenvalues of $A^*A$ are $\sigma_j^2$ by the SVD. Apply
\cref{eq:minor-scales} to both matrices.
\end{proof}

Equivalently, the positive Cauchy--Binet identities prevent leading
cancellation in the relevant sums of squared minors. Parity in
\cref{eq:gram-scales} is an arithmetic property of the group; it is not a
statement about an integer-valued valuation only.

\section{The exact field of a principal matrix root}
\label{sec:roots}

\subsection{Scalar roots are monomials times base-field units}

Every ordered abelian group is torsion free. Its divisible hull
$\Gamma_{\Q}=\Gamma\otimes_\Z\Q$ has the unique order extending that of
$\Gamma$. We use $\gamma/q$ only as an element of this hull unless its
membership in a smaller group has been established.

\begin{lemma}[The scalar root formula]\label{lem:scalarroot}
Let $a\in\FG$ be positive, let $\gamma=v(a)$, and let $q\ge2$ be an integer.
There is a unique positive $q$th root in $F_{\Gamma_{\Q}}$, and it has the form
\begin{equation}\label{eq:scalarroot}
 a^{1/q}=t^{\gamma/q}u,\qquad u\in\FG,\quad u>0.
\end{equation}
It belongs to $\FG$ if and only if $\gamma\in q\Gamma$.
\end{lemma}
\begin{proof}
Factor $a=c t^\gamma(1+h)$ with $c>0$ real and $h\in\mm\cap\FG$.
The binomial series gives
$u=c^{1/q}(1+h)^{1/q}\in\FG$, with positive constant coefficient.
This proves existence and the displayed factorization. The function
$x\mapsto x^q$ is strictly increasing for positive elements of any ordered
field, so the positive root is unique. A root in $\FG$ must have valuation
$\gamma/q\in\Gamma$, and the formula proves the converse.
\end{proof}

\begin{lemma}[Finite-index Hahn extensions]\label{lem:finite-index}
Let $\Gamma\subseteq\Lambda\subseteq\Gamma_{\Q}$ with $\Lambda/\Gamma$
finite. Then
\begin{equation}\label{eq:finite-index}
 [K_\Lambda:K_\Gamma]=[F_\Lambda:F_\Gamma]=|\Lambda/\Gamma|.
\end{equation}
A choice of coset representatives $\eta_h$ gives a basis
$(t^{\eta_h})_{h\in\Lambda/\Gamma}$ over the respective base field.
If $q\Lambda\subseteq\Gamma$, the complex extension is Galois and
\begin{equation}\label{eq:galois}
 \Gal(K_\Lambda/K_\Gamma)=\Hom(\Lambda/\Gamma,\mu_q).
\end{equation}
A character $\chi$ acts by $t^\eta\mapsto\chi(\eta+\Gamma)t^\eta$.
\end{lemma}
\begin{proof}
Split the support of a Hahn series into the finitely many cosets of $\Gamma$.
For each coset $h$, translating its support by $-\eta_h$ gives a well-ordered
subset of $\Gamma$. Thus every series has a unique expression
\[
 f=\sum_{h\in\Lambda/\Gamma}t^{\eta_h}f_h,\qquad f_h\in K_\Gamma,
\]
and the analogous real expression. Uniqueness follows because supports in
distinct cosets are disjoint. This proves the basis and degree assertion.

Multiplication by a character on each coset preserves all coefficientwise
sums and products, so defines an automorphism. A finite abelian group killed
by $q$ has exactly its own order many characters into $\mu_q$, and these
characters separate points, as is seen by decomposing it into cyclic groups.
We have therefore exhibited as many distinct automorphisms as the extension
degree. The extension is Galois with the stated group. Its residue field is
still $\C$ and its ramification index is the full degree, so it is totally
ramified.
\end{proof}

The finite-index hypothesis matters. For an infinite group extension, the
field generated by all its monomials need not contain every Hahn series
supported in the enlarged group. We will invoke \cref{lem:finite-index} only
for extensions generated by finitely many fractional exponents of bounded
denominator.

\subsection{Existence, uniqueness, and the exact extension}

\begin{theorem}[Principal-root field]\label{thm:rootfield}
Let $A\in\Mat_n(\KG)$ be Hermitian positive semidefinite, with positive
eigenvalues $\lambda_1,\ldots,\lambda_r$ counted with multiplicity, and set
$\gamma_j=v(\lambda_j)$. Define
\[
 \Lambda=\Gamma+\sum_{j=1}^r\Z\frac{\gamma_j}{q}.
\]
There is a unique positive semidefinite Hermitian matrix $B$ over
$K_{\Gamma_{\Q}}$ such that $B^q=A$. It lies in $\Mat_n(K_\Lambda)$, and
\begin{equation}\label{eq:exact-root-field}
 K_\Gamma(B)=K_\Lambda.
\end{equation}
\end{theorem}
\begin{proof}
Choose a unitary diagonalization over the base field, furnished by
\cref{thm:hermitian}:
\[
 A=U\diag(\lambda_1,\ldots,\lambda_r,0,\ldots,0)U^*.
\]
The diagonal eigenvalues are nonnegative because they are values of the
quadratic form on unit eigenvectors. By \cref{lem:scalarroot},
$b_j=\lambda_j^{1/q}=t^{\gamma_j/q}u_j$ for positive $u_j\in\FG$.
Set
\begin{equation}\label{eq:principal-root}
 B=U\diag(b_1,\ldots,b_r,0,\ldots,0)U^*.
\end{equation}
This matrix is positive semidefinite, satisfies $B^q=A$, and has entries
in $K_\Lambda$.

For uniqueness, let $C=C^*\ge0$ satisfy $C^q=A$ in the enlarged field.
The matrices $C,A$ commute, so simultaneous Hermitian diagonalization is
available there by \cref{thm:simultaneous}. On an eigenspace of $A$ with
value $\lambda$, every eigenvalue of $C$ is nonnegative and has $q$th power
$\lambda$. The ordered scalar uniqueness in \cref{lem:scalarroot} forces
it to be the chosen root, including $0$ when $\lambda=0$. Hence $C=B$.

Let $L=K_\Gamma(B)$. Since $U$ belongs to the base field, the entries of
$U^*BU$ show that every $b_j$ belongs to $L$. Dividing by the base-field
unit $u_j$ gives $t^{\gamma_j/q}\in L$. These finitely many monomials generate
all coset representatives of $\Lambda/\Gamma$ by multiplication and
inversion, together with monomials already in the base field. The group
quotient is finite, being generated by $r$ elements killed by $q$.
The basis in \cref{lem:finite-index} therefore implies $K_\Lambda\subseteq L$.
The reverse inclusion follows from \cref{eq:principal-root}.
\end{proof}

\subsection{Replacing eigenvalues by minors}

\begin{proof}[Proof of \cref{thm:headline-root}]
Order the $\gamma_j$ nondecreasingly. By \cref{eq:minor-scales},
$\Delta_k=\sum_{j=1}^k\gamma_j$. Consequently
\begin{equation}\label{eq:same-lattice}
 \Gamma+\sum_j\Z\frac{\gamma_j}q
 =\Gamma+\sum_k\Z\frac{\Delta_k}q.
\end{equation}
Indeed, each $\Delta_k/q$ is a sum of the $\gamma_j/q$, and each
$\gamma_j/q=(\Delta_j-\Delta_{j-1})/q$.

The map
\[
 \Gamma/q\Gamma\longrightarrow (q^{-1}\Gamma)/\Gamma,
 \qquad \gamma+q\Gamma\longmapsto\gamma/q+\Gamma
\]
is an isomorphism. It identifies $R_q(A)$ with
$\Gamma_{A,q}/\Gamma$. Combine \cref{thm:rootfield,lem:finite-index} to obtain
\cref{eq:field-main,eq:degree-main} and the Galois assertion. The extension
is trivial precisely when every generating class is zero, giving
\cref{eq:criterion-main}.
\end{proof}

\begin{corollary}[Exact Gram-factorization criterion]\label{cor:gram-criterion}
A Hermitian positive semidefinite matrix $A$ over $\KG$ can be written
$A=C^*C$ with $C$ over $\KG$ if and only if
\[
 \Delta_k(A)\in2\Gamma\qquad(1\le k\le\rank(A)).
\]
The condition permits $C$ to be square; it is also necessary for every
rectangular Gram representation.
\end{corollary}
\begin{proof}
Necessity is \cref{cor:gram-scales}. For sufficiency the main theorem supplies
the principal square root $B$ over $\KG$; take $C=B$.
\end{proof}

\begin{corollary}[Rational powers]\label{cor:rationalpower}
Let $A$ be positive definite and let $p\in\Z$, $q\ge2$, with
$\gcd(p,q)=1$. Then the principal power $A^{p/q}$ generates the same extension
as $A^{1/q}$. The result also holds for positive semidefinite $A$ when $p>0$.
\end{corollary}
\begin{proof}
Its nonzero diagonal spectral entries are $b_j^p$. If $ap+bq=1$ for integers
$a,b$, then $b_j=(b_j^p)^a\lambda_j^b$, so adjoining these entries is
field-theoretically equivalent to adjoining the $b_j$. Zero eigenvalues
are omitted; negative powers require invertibility.
\end{proof}

\subsection{Prime and composite denominators}

When $q$ is prime, $\Gamma/q\Gamma$ is a vector space over $\mathbf F_q$,
and
\[
 [K_\Gamma(A^{1/q}):K_\Gamma]
 =q^{\dim_{\mathbf F_q}\Span\{\Delta_k+q\Gamma:1\le k\le r\}}.
\]
For composite $q$ the exact invariant is the finite subgroup, not the rank of
a vector space. For example, a single class of order $2$ modulo $4\Gamma$
contributes a quadratic extension, not a quartic one.

If $\Gamma=\Z^d$ with any fixed compatible order and the $\Delta_k$ are
integer vectors, let
\[
 L=q\Z^d+\sum_{k=1}^r\Z\Delta_k\ \subseteq\ \Z^d.
\]
Then
\begin{equation}\label{eq:lattice-degree}
 |R_q(A)|=\frac{q^d}{[\Z^d:L]}.
\end{equation}
The index in the denominator is computed by Smith normal form of the
integer matrix with columns $qI_d,\Delta_1,\ldots,\Delta_r$. Thus the root
field degree is a finite integer-lattice computation after the minor
valuations have been determined exactly. No discreteness in the \emph{order}
of $\Gamma$ is needed for this algebraic computation.

\section{The trace is a primitive generator}
\label{sec:trace}

\subsection{A positive-coset lemma}

\begin{lemma}[Positive radical sums retain every generator]\label{lem:positive-trace}
Let $a_1,\ldots,a_r\in\FG$ be positive, let $q\ge2$, and write
$b_j=a_j^{1/q}=t^{\gamma_j/q}u_j$ as in \cref{lem:scalarroot}. Put
$\Lambda=\Gamma+\sum_j\Z\gamma_j/q$. If $w_j\in\FG$ are strictly positive,
then
\begin{equation}\label{eq:positive-sum}
 K_\Gamma\left(\sum_{j=1}^r w_jb_j\right)=K_\Lambda.
\end{equation}
\end{lemma}
\begin{proof}
Let $G=\Lambda/\Gamma$ and let $h_j=\gamma_j/q+\Gamma$. These classes
generate $G$. Choose a representative $\eta_h$ for each class occurring
among the $h_j$. The sum $T=\sum_jw_jb_j$ has the coset decomposition
\[
 T=\sum_h t^{\eta_h}s_h,\qquad
 s_h=\sum_{j:h_j=h}w_j\,t^{\gamma_j/q-\eta_h}u_j\in\FG.
\]
Every summand in $s_h$ is positive in the ordered field $\FG$; hence
$s_h>0$ and in particular $s_h\ne0$ for each occurring class.

By \cref{lem:finite-index}, all automorphisms of $K_\Lambda/K_\Gamma$ are
characters $\chi:G\to\mu_q$. If such an automorphism fixes $T$, then
\[
 0=\chi(T)-T=\sum_h(\chi(h)-1)t^{\eta_h}s_h.
\]
Distinct cosets have disjoint supports, so every coefficient $\chi(h)-1$
must be zero. Since the occurring $h_j$ generate $G$, the character is
trivial. Thus $T$ has trivial stabilizer in the Galois group. The fundamental
correspondence, or the orbit-degree formula, gives
$[K_\Gamma(T):K_\Gamma]=|G|$, proving the equality.
\end{proof}

\begin{proof}[Proof of \cref{thm:headline-trace}]
For the principal root in \cref{eq:principal-root}, trace invariance gives
$\tr(B)=\sum_{j=1}^r\lambda_j^{1/q}$. Apply
\cref{lem:positive-trace} with $w_j=1$ and use
\cref{thm:rootfield,eq:same-lattice}. The weighted statement is the same
lemma with arbitrary strictly positive $w_j$.
\end{proof}

The proof identifies the mechanism rather than merely invoking an abstract
primitive-element theorem. In a general finite extension, most linear
combinations may be primitive, but special choices can fail. Here the
\emph{specific} unweighted trace always succeeds because positivity makes
its coset support a generating set.

\begin{corollary}[The corresponding real field]\label{cor:realtrace}
With $T=\tr(A^{1/q})$ and $\Lambda=\Gamma_{A,q}$,
\[
 F_\Gamma(T)=F_\Lambda.
\]
\end{corollary}
\begin{proof}
Both fields are ordered, and $T\in F_\Lambda$. Adjoining $i$ has degree $2$
over every one of these real fields. Applying the tower formula to
$K_\Gamma(T)=K_\Lambda$ shows that
$[F_\Gamma(T):F_\Gamma]=[K_\Lambda:K_\Gamma]=|\Lambda/\Gamma|$.
This is also $[F_\Lambda:F_\Gamma]$ by \cref{lem:finite-index}, so the
inclusion is equality.
\end{proof}

\subsection{An explicit irreducible orbit polynomial}

Let $b_j=\lambda_j^{1/q}$, $G=\Gamma_{A,q}/\Gamma$, and
$h_j=v(b_j)+\Gamma$. Then the minimal polynomial of the trace over $\KG$ is
\begin{equation}\label{eq:orbit-poly}
 \mathcal P_{A,q}(X)=
 \prod_{\chi\in\Hom(G,\mu_q)}
 \left(X-\sum_{j=1}^r\chi(h_j)b_j\right).
\end{equation}
The coefficients belong to $\KG$ because the Galois group permutes the
factors. The factors are distinct by \cref{lem:positive-trace}, so the
polynomial is irreducible and has degree $|R_q(A)|$. It is a formula inside
a finite extension, not necessarily a fast coefficient-extraction algorithm.

\begin{example}[A biquadratic field encoded in one positive trace]
Let $\Gamma=\Z\delta\oplus\Z\Omega$, with $0<n\delta<\Omega$ for all
positive integers $n$, and let
\[
 A=\diag(t^\delta,t^\Omega),\qquad q=2.
\]
The minor valuations are $\Delta_1=\delta$ and
$\Delta_2=\delta+\Omega$. Their classes span a subgroup of order $4$ in
$\Gamma/2\Gamma$. Thus $B=A^{1/2}$ generates a biquadratic extension, and
so does
$T=t^{\delta/2}+t^{\Omega/2}$. Writing $t=t^\delta$, $u=t^\Omega$, the
minimal polynomial is
\begin{equation}\label{eq:biquadratic}
 X^4-2(t+u)X^2+(t-u)^2.
\end{equation}
The degree $4$ is certified by the two independent exponent classes, not
by an unproved irreducibility guess for the displayed polynomial.
\end{example}

\begin{warning}[Positivity cannot be dropped]
For $A=t^\delta I_2$ with $\delta\notin2\Gamma$, the matrix
$C=\diag(t^{\delta/2},-t^{\delta/2})$ is a Hermitian square root and generates
a quadratic extension, but $\tr(C)=0$. It is not the positive root.
The main trace theorem concerns the principal positive root and strictly
positive weights, not arbitrary branches or signed sums.
\end{warning}

\subsection{Several positive roots at once}

\begin{corollary}[Joint ramification and a joint primitive element]\label{cor:joint}
Let $A_1,\ldots,A_m$ be positive semidefinite Hermitian matrices over
$\KG$, of possibly different sizes, and let $q_\ell\ge2$.
Put
\[
 \Lambda=\Gamma+
  \sum_{\ell=1}^m\sum_{k=1}^{\rank(A_\ell)}
       \Z\frac{\Delta_k(A_\ell)}{q_\ell}.
\]
Then the field generated by all entries of all $A_\ell^{1/q_\ell}$ is
$K_\Lambda$. Its degree is $|\Lambda/\Gamma|$, and the single scalar
\[
 T=\sum_{\ell=1}^m\tr(A_\ell^{1/q_\ell})
\]
is a primitive generator. No commutation assumption on the $A_\ell$ is
necessary.
\end{corollary}
\begin{proof}
The compositum of the individual monomial extensions contains the monomials
generating $\Lambda$. The quotient is finite and killed by
$Q=\lcm(q_1,\ldots,q_m)$, so \cref{lem:finite-index} identifies the compositum
with $K_\Lambda$. Decompose the finitely many positive scalar spectral roots
in the traces into cosets of $\Gamma$. Every occurring coefficient sum is
positive, and the occurring classes generate $\Lambda/\Gamma$. The character
argument of \cref{lem:positive-trace}, now with $\mu_Q$, proves primitivity.
\end{proof}

\section{Consequences inside the surreal and surcomplex fields}
\label{sec:surreal}

\subsection{Localization to the actual input support}

The normal-form theorem for surreal numbers expresses each surreal as a
set-supported real linear combination of monomials $\omega^a$, with reverse
well-ordered exponents \cite{Gonshor}. Replacing $a$ by $-\gamma$ gives our
increasing Hahn convention $t^\gamma=\omega^{-\gamma}$.
The algebraic compatibility of these normal forms identifies the Hahn field
over any set-sized ordered subgroup $H$ of the additive surreal group with
a subfield of $\No$; complexification embeds $K_H$ into $\SC$.

For a finite surcomplex matrix $A$, take the union of the supports of the
real and imaginary normal forms of all its entries, negate the exponents,
and let $H$ be the additive subgroup they generate. This is a set, because
a subgroup generated by a set consists of finite integer combinations.
Every entry of $A$ lies in $K_H$. No global class-sized Hahn sum is required.

\begin{corollary}[Exact-support surcomplex spectral calculus]\label{cor:surcomplex}
A Hermitian surcomplex matrix has real surreal eigenvalues and a unitary
eigenbasis whose normal-form exponents stay in the additive group generated
by the input exponents. The same support-group assertion holds for normal
diagonalization and for the SVD of an arbitrary finite surcomplex matrix.

For a positive semidefinite Hermitian matrix and its exact support group $H$,
the principal $q$th root has the exact entry field
\[
 K_H(A^{1/q})
   =K_{H+\sum_k\Z\Delta_k(A)/q}
   =K_H\bigl(\tr(A^{1/q})\bigr).
\]
The extension degree is the order of the subgroup generated by the
$\Delta_k(A)$ modulo $qH$.
\end{corollary}
\begin{proof}
Apply the preceding theorems in the set-sized field $K_H$ and then use the
normal-form embedding. The positive branch agrees with the ambient surreal
positive root by ordered-field uniqueness and the uniqueness proof for the
matrix root.
\end{proof}

The ambient surreal exponent group is divisible, so the full class field has
room for all the displayed fractional scales. The content of the corollary
is not a failure of positive roots to exist in $\No$. It is their exact
\emph{relative cost} over the smallest Hahn workspace containing the data.
This cost is visible before one replaces $H$ by its divisible hull.

\subsection{Transport between ordered scale systems}

\begin{proposition}[Order-preserving transport]\label{prop:transport}
Let $\phi:\Gamma\hookrightarrow\Lambda$ be an injective order-preserving
group homomorphism. The map
\[
 \phi_*\left(\sum_\gamma a_\gamma t^\gamma\right)
       =\sum_\gamma a_\gamma s^{\phi(\gamma)}
\]
is an injective star-preserving field homomorphism. It preserves strongly
summable families, the support certificates of the splitting construction,
and each chosen spectral or singular-value decomposition.
\end{proposition}
\begin{proof}
An order embedding sends a well-ordered support to a well-ordered support.
It preserves and reflects equality of sums of exponents, so the finite
coefficient formulas for strong sums and products are unchanged. It also
preserves leading coefficients and minima. Thus the trace subtraction,
normalization, formal projectors, polar intertwiner, and all selected ordinary
constant eigenbases transport through every node of the finite tree.
\end{proof}

This permits one symbolic support construction to be interpreted at different
surreal scales having the same ordered-group relations. It does not claim
that arbitrary numerical substitutions, especially ones that identify or
reverse exponents, preserve Hahn summability.

\begin{warning}[Relative degrees can decrease in a larger workspace]
Under an embedding into a larger group, an exponent not divisible by $q$ in
the source may become divisible in the target. For instance,
$\Z\hookrightarrow\Q$ makes every positive monomial's $q$th root available.
The spectral identities transport, but the relative ramification group must
be recomputed in the new base group. Degree is preserved under an isomorphism
onto the image group, not under arbitrary enlargement of the base workspace.
\end{warning}

\section{Examples, sharpness, and category boundaries}
\label{sec:examples}

\subsection{No new fractions, but new integer combinations}

For $\Gamma=\Z$ and $t=t^1$, consider
\begin{equation}\label{eq:t3matrix}
 A=\begin{pmatrix}t^3&t^4\\t^4&0\end{pmatrix}.
\end{equation}
Its exact eigenvalues are
\[
 \lambda_\pm=\frac{t^3}{2}\left(1\pm\sqrt{1+4t^2}\right).
\]
The first terms are
\begin{align*}
 \lambda_+&=t^3+t^5-t^7+2t^9-5t^{11}+14t^{13}-42t^{15}+\cdots,\\
 \lambda_-&=-t^5+t^7-2t^9+5t^{11}-14t^{13}+42t^{15}+\cdots.
\end{align*}
Define
\[
 w=\frac{\sqrt{1+4t^2}-1}{2t}
      =t-t^3+2t^5-5t^7+14t^9-\cdots.
\]
Then the orthogonal matrix
\[
 U=(1+w^2)^{-1/2}
      \begin{pmatrix}1&-w\\w&1\end{pmatrix}
\]
satisfies $U^*AU=\diag(\lambda_+,\lambda_-)$. All exponents are integers.
However, exponent $5$ is not in the nonnegative monoid generated by the raw
input exponents $3$ and $4$. Therefore a theorem claiming that the eigenvalue
supports always remain in the \emph{raw input monoid} would be false.
The correct invariant is the input \emph{subgroup}, supplemented by the
positive monoids of normalized perturbations at each node.

\subsection{An eigenvalue beyond all ordinary perturbation orders}

In the rank-two group of \cref{ex:nonconvergence}, abbreviate
$t=t^\delta$ and $u=t^\Omega$ and consider
\[
 A=\begin{pmatrix}0&u\\u&t\end{pmatrix}.
\]
Since $v(u/t)=\Omega-\delta>0$, its smaller eigenvalue is
\begin{align*}
 \lambda_-&=\frac t2\left(1-\sqrt{1+4u^2/t^2}\right)\\
 &=-u^2t^{-1}+u^4t^{-3}-2u^6t^{-5}+5u^8t^{-7}-\cdots.
\end{align*}
The valuation is $2\Omega-\delta$, which exceeds every finite multiple of
$\delta$. The support remains in $\Z\delta+\Z\Omega$, and the expansion is
certified by the positive generator $2(\Omega-\delta)$. A finite-order
expansion measured only in powers of $t$ does not see this eigenvalue.

\subsection{A positive matrix whose square root genuinely ramifies}

Take
\begin{equation}\label{eq:positive2}
 A=\begin{pmatrix}1+t&1\\1&1\end{pmatrix}\in\Mat_2(\C((t))).
\end{equation}
It is positive definite: for $x=(x_1,x_2)$,
\[
 x^*Ax=|x_1+x_2|^2+t|x_1|^2>0\quad(x\ne0).
\]
The minor minima are $\Delta_1=0$, $\Delta_2=v(t)=1$. Therefore the
principal square root generates exactly $\C((t^{\frac12\Z}))$, a quadratic
extension of $\C((t^\Z))$. An explicit formula is
\begin{equation}\label{eq:positive2root}
 B=\frac{A+\sqrt t\,I}{\sqrt{2+t+2\sqrt t}}.
\end{equation}
The denominator has a positive constant leading term, so its square root
introduces no exponent beyond $\tfrac12\Z$. Cayley--Hamilton gives
$B^2=A$. The positive numerator and positive denominator show $B\ge0$.
Moreover
\[
 T=\tr(B)=\sqrt{2+t+2\sqrt t},\qquad
 \sqrt t=\frac{T^2-2-t}{2},
\]
which makes the primitive-trace phenomenon directly visible.

\subsection{A Gram witness can use a larger group than its product}

Let $t=t^\delta$, $u=t^\Omega$ in the same rank-two group and put
\[
 C=\begin{pmatrix}t&1\\0&u\end{pmatrix},\qquad
 A=C^*C=\begin{pmatrix}t^2&t\\t&1+u^2\end{pmatrix}.
\]
Relative to $\Gamma=\Z\delta+\Z\Omega$, the minor valuations are
$\Delta_1=0$, $\Delta_2=2\delta+2\Omega$. The square root requires no
extension; the fourth root generates an extension of degree $2$.

But the \emph{exact input support group of $A$ alone} is
$H=\Z\delta+2\Z\Omega$. Relative to $H$, the value
$2\delta+2\Omega$ is not in $2H$, so $A^{1/2}$ generates a quadratic
extension and recovers the missing scale $u$. This does not contradict
\cref{cor:gramroot}: the displayed Gram witness $C$ already uses that scale.
It illustrates why statements must name both the input and the base group.

\subsection{The degree bound is attained in every rank}

Let $\Gamma=\Z^r$ with lexicographic order and standard basis
$e_1,\ldots,e_r$. Each $e_j$ is positive. For
\[
 A=\diag(t^{e_1},\ldots,t^{e_r}),
\]
the nonzero eigenvalue valuations generate all of $\Gamma/q\Gamma$.
Their partial sums, which are the determinantal values after ordering,
generate the same subgroup. Hence
\[
 [K_\Gamma(A^{1/q}):K_\Gamma]=q^r.
\]
This sharpness example rules out a degree bound depending only on $q$ and
not on the matrix rank. Its trace is already a primitive element of the
maximal-degree extension.

\subsection{Why this is not a multivariable analytic eigenbasis theorem}

Consider the real symmetric formal matrix
\begin{equation}\label{eq:crossing}
 H(x,y)=\begin{pmatrix}x&y\\y&-x\end{pmatrix}.
\end{equation}
Its characteristic equation is $\lambda^2=x^2+y^2$. There is no solution
$\lambda\in\C[[x,y]]$: its lowest homogeneous term would have to be
$ax+by$ with $a^2=b^2=1$ and $2ab=0$, which is impossible.
Thus a general eigenvalue or eigenbasis theorem in that formal power-series
ring is false. This is the kind of distinction relevant to the
multiparameter literature \cite{PR,DLLZ}.

After the ordered Hahn specialization $x=t^\alpha$, $y=t^\beta$ with
$\beta>\alpha$, however,
\[
 \sqrt{x^2+y^2}
  =t^\alpha\sum_{m\ge0}\binom{1/2}{m}t^{2m(\beta-\alpha)}
\]
is a legitimate Hahn series in the original group. It uses the ordering and
the allowed division by the dominant monomial. The displayed support may
correspond to negative powers of $x$ in the original two-variable expression.
It therefore does not give an element of $\C[[x,y]]$, a common analytic germ
through the crossing, or a global analytic choice of eigenvectors. Our result
resolves the totally ordered scale problem, not those different problems.

\section{Computation and a verification contract}
\label{sec:computation}

\subsection{What must be supplied by a symbolic representation}

The spectral construction needs exact arithmetic and conjugation, exact
zero tests, leading exponents and coefficients for nonzero expressions,
ordinary constant Hermitian diagonalization, and a representation of
strongly summable substitutions with their support certificates. An exact
representation need not provide any of these effectively merely because it
can denote a surreal number.

The determinantal ramification computation needs less: exact valuations of
finitely many minors and a way to compute the finite subgroup generated by
specified classes modulo $q\Gamma$. In a finite integer lattice,
\cref{eq:lattice-degree} gives a concrete algorithm. For an arbitrary abstract
ordered group, membership in $q\Gamma$ is additional group data and need
not be decidable from a black-box order comparison.

A practical implementation should keep two outputs separate. A matrix
identity such as $U^*AU=D$ is an algebraic certificate. A support certificate
licenses the Hahn expressions occurring in that identity. Finite truncations
can test an identity to a declared order, but they do not certify exact zero,
all coefficients, or the well-ordering of arbitrary supports.

\subsection{Included exact checks}

The companion script \texttt{code/verify.py} performs rational and symbolic
checks, with no floating-point sampling. It exercises the resolvent-projector
construction and its polar intertwiner through formal degree six for a
three-dimensional Hermitian example with a repeated leading eigenvalue;
it tests the displayed scalar and two-by-two identities; it checks
minor-valued profiles under rational unitary changes of basis; and it tests
finite quotient computations and character separation in finite integer
lattices. The executed suite passed \textbf{807 exact scalar assertions}: 378 projector
and polar-series coefficients, 55 block and low-degree terms, 29 displayed
identities, 144 finite ramification-group checks, 96 character checks, and
105 rational-matrix checks. It includes 48 finite quotient cases and six
rational matrix profiles. The run used Python 3.13.5 and SymPy 1.14.0.
Counts include individual zero coefficients and are not counts of independent
theorems. Full results are recorded in \texttt{data/verification.json} and
\texttt{data/verification.txt}.

These checks are useful for detecting sign, conjugation, denominator, and
valuation-convention errors. They are not a machine proof of the Hahn support
lemma, the arbitrary-rank spectral theorem, or the primitive-trace theorem.
The general arguments are the mathematical proofs in this article. No Lean
formalization of these new results is included or claimed.

\subsection{A plausible formalization order}

The existing repository already has a Hahn summability layer and extensive
finite polynomial algebra \cite{RepoMain}. A continuation toward the present
results could first formalize positive sum-of-squares valuation and integral
unitary invariance. Next come finite-index Hahn coset decompositions, the
monomial character action, and the positive-coset primitive-element argument.
These are set-sized statements and do not depend on completing the full
surreal field construction.

The spectral construction needs an interface for finite-variable formal
series evaluated at Hahn infinitesimals, then the universal separated-block
projectors and their polar intertwiner. The ordinary complex contour
justification in \cref{thm:split} could be replaced in a proof assistant by
formal resolvent identities or a formal Hensel block-lifting theorem.
Dimension induction then supplies the finite splitting tree. Finally,
normal-form embeddings transport the results to whichever universe-sized
surreal carrier is available. This is a proposed dependency order, not a
report of code that has been compiled.

\section{Conclusions and remaining questions}
\label{sec:conclusion}

The article proves a distinction that is obscured by working only in the
full surreal class field. Hermitian eigenvalues, normal eigenvectors, and
singular-value decompositions need no fractional enlargement of the actual
Hahn support group. Positive matrix roots do sometimes need such an
enlargement, and the required field is exactly determined by valuations of
minors. A single scalar---the trace of the principal root---generates that
whole field.

The proofs apply at arbitrary ordered-group rank. Their infinite operations
are justified by word-finite Hahn support rather than by an assumed
cofinal sequence of scales. The finite bound on splitting nodes is a bound
on spectral dimension reduction, not on the number of coefficients to be
read. The sharp examples distinguish subgroup closure from raw monoid
closure, a matrix from a larger-support Gram witness, and totally ordered
Hahn specialization from multiparameter formal analyticity.

Three further research directions are particularly concrete. One is to
replace the finite tree of set-valued supports by finite effective support
grammars for useful transseries subclasses. Another is to characterize
smaller non-Hahn subfields that are closed under the separated-block
construction, so that the full Hahn field can be replaced by a more
computationally economical domain. A third is to investigate which other
positive spectral functionals, beyond rational powers and their strictly
positive sums, furnish similarly exact primitive generators. These are
questions suggested by the results, not advertised as previously published
open problems.

The novelty assessment remains deliberately narrow. The earlier
power-series spectral theorems are real predecessors, and the ordinary
algebraic ingredients are standard. The contribution offered for further
expert examination is the exact arbitrary-rank support-descent and
determinantal/primitive-trace synthesis. Its correctness is supported here
by explicit proofs and finite checks; its priority and a full formal
verification remain unestablished.

\appendix
\section{Dependency and hypothesis audit}
\label{app:audit}

\begin{longtable}{@{}L{0.25\textwidth}L{0.35\textwidth}L{0.34\textwidth}@{}}
\toprule
Result & Essential inputs & What is not assumed \\
\midrule\endhead
Positive norm square & Real coefficients, leading order, positive support binomial evaluation & Square roots of all positive field elements \\
\addlinespace
Separated-block lift & Distinct ordinary residue eigenvalues, finite-variable formal identities, positive Hahn support & Sequential valuation convergence or a surreal contour integral \\
\addlinespace
Hermitian spectral descent & Trace subtraction, constant complex spectral theorem, block lift, induction on dimension & Divisibility of $\Gamma$, real closedness of $F_\Gamma$, finite rank \\
\addlinespace
Normal and simultaneous descent & Hermitian descent, commutant argument, finite-dimensional refinement & A pre-existing algebraic closure of $K_\Gamma$ or an infinite product of unitaries \\
\addlinespace
SVD & Hermitian Gram diagonalization, norm-square roots, finite Gram--Schmidt & Square roots of arbitrary positive eigenvalues of arbitrary positive matrices \\
\addlinespace
Minor scales & Cauchy--Binet, integral inverse transformations, finite minima & A discrete or Noetherian valuation ring \\
\addlinespace
Exact root field & Scalar binomial roots, spectral descent, finite-index coset basis & A generic-eigenvalue assumption, nonzero determinant, or simple spectrum \\
\addlinespace
Primitive trace & Complex roots of unity, coset character action, strictly positive coefficients & Generic weights or independent eigenvalue classes \\
\addlinespace
Surcomplex interpretation & Set-supported surreal normal forms, set-sized input subgroup & Proper-class vector spaces, proper-class sums, or same-universe smallness of all surreals \\
\bottomrule
\end{longtable}

\clearpage
\section{A compact exact-operation description}
\label{app:algorithm}

For a Hermitian matrix, the recursive routine is the following. It is
pseudocode with exact Hahn operations, not executable code for arbitrary
surreal representations.
\begin{quote}\ttfamily\small
Spectral(A):\\
\hspace*{1em}n := number of rows of A\\
\hspace*{1em}mu := trace(A) / n; C := A - mu*I\\
\hspace*{1em}if C = 0: return (I, mu*I, scalar-leaf)\\
\hspace*{1em}alpha := minimum valuation of the entries of C\\
\hspace*{1em}C0 := residue(t\string^(-alpha)*C)\\
\hspace*{1em}(V,D,blocks) := ordinary spectral decomposition of C0\\
\hspace*{1em}E := t\string^(-alpha)*V* C V - D\quad [V* means adjoint]\\
\hspace*{1em}(P,W,certificate) := separated-block lift of (D,E)\\
\hspace*{1em}Ablocks := W* V* A V W\\
\hspace*{1em}recurse on every diagonal block of Ablocks\\
\hspace*{1em}return the product basis, diagonal, and certificate tree
\end{quote}
The products in this pseudocode use $V^*$ and $W^*$ as adjoints, not
programming-language multiplication symbols. The exact-operation meaning is
fixed by the displayed mathematical formulas in \cref{sec:spectral}.

For ramification alone, enumerate minors to compute $\Delta_1,\ldots,
\Delta_r$, form their classes modulo $q\Gamma$, and compute the finite
subgroup they generate. Return that group, its order, and the extension
$\Gamma+\sum_k\Z\Delta_k/q$. No spectral decomposition is required to
produce this arithmetic answer once the positivity hypothesis and exact
minor valuations are known.

\begin{thebibliography}{99}
\bibitem{Adkins}
W.~A.~Adkins,
\emph{Normal matrices over hermitian discrete valuation rings},
Linear Algebra and its Applications \textbf{157} (1991), 165--174.
\url{https://doi.org/10.1016/0024-3795(91)90111-9}.
The author's institutional record was consulted for the stated scope:
\url{https://repository.lsu.edu/mathematics_pubs/10/}.

\bibitem{KO}
H.~A.~Keller and H.~Ochsenius A.,
\emph{A spectral theorem for matrices over fields of power series},
Annales math\'ematiques Blaise Pascal \textbf{2}, no.~1 (1995), 169--179.
\url{https://ambp.centre-mersenne.org/item/AMBP_1995__2_1_169_0/}.
The full article, including the groups in Theorems 2 and 3, was consulted.

\bibitem{PR}
A.~Parusi\'nski and G.~Rond,
\emph{Multiparameter perturbation theory of matrices and linear operators},
Transactions of the American Mathematical Society \textbf{373} (2020),
2933--2948.
\url{https://doi.org/10.1090/tran/8061};
\url{https://arxiv.org/abs/1807.04242}.
The abstract and its theorem descriptions in the authors' subsequent
literature were consulted; the block-lifting argument used here is proved
independently in \cref{thm:split}.

\bibitem{DLLZ}
Z.~Dai, H.~Liang, J.~Lu, and L.~Zhi,
\emph{An Algorithm for Diagonalizing Matrices of Formal Power Series},
arXiv:2602.08313v1, February 9, 2026.
\url{https://arxiv.org/html/2602.08313v1}.
A preprint, not described here as a refereed publication.

\bibitem{Neumann}
B.~H.~Neumann,
\emph{On ordered division rings},
Transactions of the American Mathematical Society \textbf{66} (1949),
202--252.
\url{https://www.jstor.org/stable/1990552}.
The support lemma is reproved in \cref{lem:neumann}; the bibliographic record,
not a fully accessible transcription of the original article, was consulted.

\bibitem{Higman}
G.~Higman,
\emph{Ordering by divisibility in abstract algebras},
Proceedings of the London Mathematical Society (3) \textbf{2} (1952), 326--336.
\url{https://doi.org/10.1112/plms/s3-2.1.326}.
The needed ordered-word special case is proved in \cref{lem:words}.

\bibitem{Gonshor}
H.~Gonshor,
\emph{An Introduction to the Theory of Surreal Numbers},
London Mathematical Society Lecture Note Series 110,
Cambridge University Press, 1986.
\url{https://www.cambridge.org/core/books/an-introduction-to-the-theory-of-surreal-numbers/312AE504A3E88E804054BFB390446374}.
The normal-form and ordered-field framework is an imported foundational input,
not reconstructed in this article.

\bibitem{RepoMain}
V.~Reshetnikov, \emph{Surreal}, repository README and documentation map,
snapshot \Commit, inspected September 21, 2026.
\url{https://github.com/VladimirReshetnikov/Surreal/blob/aa846271b4dcae2c055b216126a87210292ec19b/README.md};
\url{https://github.com/VladimirReshetnikov/Surreal/blob/aa846271b4dcae2c055b216126a87210292ec19b/docs/README.md}.

\bibitem{RepoSpectral}
\emph{Surcomplex Spectral Theory: Singular Values, Infinitesimal Rank,
and Multiscale Stability}, research package in \cite{RepoMain},
\texttt{docs/surcomplex/spectral-theory/}.
\url{https://github.com/VladimirReshetnikov/Surreal/blob/aa846271b4dcae2c055b216126a87210292ec19b/docs/surcomplex/spectral-theory/README.md}.
The full README was consulted. The package identifies itself as an
AI-assisted, unrefereed report and makes no priority claim.

\bibitem{RepoCAS}
\emph{Computer algebra}, research package in \cite{RepoMain},
\nolinkurl{docs/foundations-and-computation/computer-algebra/}.
The distinctions between exact denotation, coefficient access, decidable
equality, and approximation were consulted through the repository's detailed
documentation map, not through a full audit of the package's implementations.
\end{thebibliography}
\end{document}
